\section{Conclusion and Outlook}\label{conclusion}

We have presented an IaC framework that brings Terraform/OpenTofu into CERN's existing provisioning and CM ecosystem. It adds a declarative integration layer that holds infrastructure intent, provisioning, and service registration under one version-controlled roof. With custom providers for the in-house services and pipelines that gate every change, it cuts boundary drift, makes change traceable, and reduces how much undocumented operational knowledge a team needs to keep a service alive. That matters most in BC/DR and cross-environment replication.

Real deployments and operator feedback both point the same way: the framework makes operations more predictable, and it does so without disturbing the CM workflows people already rely on. Separating provisioning intent, execution, and ownership of runtime configuration works for infrastructure that is large, heterogeneous, and long-lived, which describes CERN's rather well.

There is more to do. Tying the framework to GitOps-style reconciliation controllers would allow continuous drift correction for selected resources, rather than waiting for an operator to re-apply. Adoption will be gradual: teams who know the existing workflows and have never used Terraform need time and support, and the migration path has to be incremental. Finally, making Policy-as-Code enforcement a design-time constraint rather than an optional check would put security and compliance guarantees on a firmer footing across every managed resource.
